\begin{mistake}{2026-05-26}{05}

\begin{problemcontent}
设 \(n\) 为自然数，求反常积分：
\[
I_n = \int_0^{+\infty} \frac{\mathrm{d}x}{(x^2+1)^n}
\]
\end{problemcontent}

\begin{solutioncontent}
\textbf{法一：三角代换 + 华里士公式}
令 \(x = \tan t\)，\(t \in [0, \frac{\pi}{2})\)，则 \(\mathrm{d}x = \sec^2 t \mathrm{d}t\)。代入原积分得：
\[
\begin{aligned}
I_n &= \int_0^{\frac{\pi}{2}} \frac{1}{(\tan^2 t + 1)^n} \cdot \sec^2 t \mathrm{d}t \\
    &= \int_0^{\frac{\pi}{2}} \frac{\sec^2 t}{\sec^{2n} t} \mathrm{d}t \\
    &= \int_0^{\frac{\pi}{2}} \cos^{2n-2} t \mathrm{d}t
\end{aligned}
\]
利用华里士公式：
当 \(n=1\) 时，\(I_1 = \int_0^{\frac{\pi}{2}} 1 \mathrm{d}t = \frac{\pi}{2}\)；
当 \(n \ge 2\) 时，\(2n-2\) 为正偶数，故 \(I_n = \frac{(2n-3)!!}{(2n-2)!!}\frac{\pi}{2}\)。

\textbf{法二：分部积分递推法}
对 \(I_{n-1}\)（\(n \ge 2\)）使用分部积分：
\[
\begin{aligned}
I_{n-1} &= \int_0^{+\infty} \frac{1}{(x^2+1)^{n-1}} \mathrm{d}x \\
&= \left[ \frac{x}{(x^2+1)^{n-1}} \right]_0^{+\infty} - \int_0^{+\infty} x \mathrm{d}\left( \frac{1}{(x^2+1)^{n-1}} \right) \\
&= 0 + 2(n-1) \int_0^{+\infty} \frac{x^2}{(x^2+1)^n} \mathrm{d}x
\end{aligned}
\]
利用“加一减一”法凑分母结构：
\[
\begin{aligned}
\int_0^{+\infty} \frac{x^2}{(x^2+1)^n} \mathrm{d}x &= \int_0^{+\infty} \frac{(x^2+1)-1}{(x^2+1)^n} \mathrm{d}x \\
&= I_{n-1} - I_n
\end{aligned}
\]
代回原式得 \(I_{n-1} = 2(n-1)(I_{n-1} - I_n)\)。
整理得递推公式：
\[
I_n = \frac{2n-3}{2n-2} I_{n-1} \quad (n \ge 2)
\]
又因初始项 \(I_1 = \int_0^{+\infty} \frac{1}{x^2+1} \mathrm{d}x = \frac{\pi}{2}\)，不断递推可得 \(I_n = \frac{(2n-3)!!}{(2n-2)!!}\frac{\pi}{2}\)。
\end{solutioncontent}

\mistakeknowledge{
\begin{itemize}
  \item \textbf{核心公式}：华里士（点火）公式偶数次幂情形：\(\int_0^{\frac{\pi}{2}} \cos^{2k} x \mathrm{d}x = \frac{(2k-1)!!}{(2k)!!}\frac{\pi}{2}\)（末尾必乘 \(\frac{\pi}{2}\)）。
  \item \textbf{解题技巧}：被积函数含 \((x^2+a^2)^n\) 时首选三角代换 \(x=a\tan t\)；有理函数积分求递推式常使用“分部积分 + 加一减一拆项法”。
  \item \textbf{易错点}：做三角代换时极易漏乘微分项 \(\mathrm{d}x = \sec^2 t \mathrm{d}t\)，导致整体降幂错误。
\end{itemize}}

\end{mistake}
